% !LW recipe = latexmk (xelatex)
\documentclass[12pt,a4paper,UTF8]{ctexart}
\usepackage{SJTUReport}
\usepackage{booktabs} % 三线表
\usepackage{amsmath}
\usepackage{amssymb}

%%-------------------------------正文开始---------------------------%%
\begin{document}

% %-----------------------封面--------------------%%
\cover

%------------------摘要-------------%%
\newpage
\begin{abstract}

面向高速公路长途出行场景，本文研究沿线换电站网络中预约与随机两类异质用户的实时运营优化问题。长途电动汽车在一次行程中可能需要多次换电，每次换电在取走满电电池的同时退回一块低电量电池，使用户的换电站序列选择与沿线各站的电池库存形成时空耦合。为此，本文构建模型预测控制（MPC）与强化学习（RL）相结合的分层协同框架：MPC 在滚动预测域内联合优化预约用户的剩余换电站序列与站内连续服务事件安排，RL actor 为各站逐槽电池生成请求充电功率，critic 则以终端 SOC 边际价值和边界未决交付价值刻画预测域之外的长期影响。模型以连续时间刻画到站、等待、超时、充电完成与瞬时换电事件，允许未决等待请求跨滚动边界携带，并按真实发生时刻对服务、违约与路径发布结算一次。

\end{abstract}
\newpage

\thispagestyle{empty} % 首页不显示页码

%%--------------------------目录页------------------------%%
\newpage
\tableofcontents

%%------------------------正文页从这里开始-------------------%
\newpage

%可选择这里也放一个标题
\begin{center}
   \title{ \Huge \textbf{{考虑异质用户的高速公路换电站运营 MPC-RL 分层协同优化}}}
\end{center}


\section{Introduction}
随着电动汽车在高速公路中长距离出行场景中的普及，沿线补能设施的运营能力逐渐成为影响用户出行可靠性与系统运行效率的关键因素。相比传统充电方式，换电模式可以缩短车辆补能时间。在高速公路场景下，长途出行的电动汽车往往需要沿行程多次换电：车辆沿道路单向行驶且不能折返，前一次换电的站点选择直接改变后续可达站点集合与剩余换电站序列，各次换电决策因此相互关联。

与此同时，不同车辆到达时的电池荷电状态（SOC）各不相同：SOC 通过影响剩余续航里程，决定车辆首个可达换电站以及整条换电站序列；而每次换电在换出满电电池的同时退回一块低 SOC 电池，退回电池须经重新充电后才能再次服务。因此，一位用户的站点选择不仅影响自身后续行程，还会改变相关站点未来的电池库存状态，使用户路径与多站点电池库存产生时空耦合。高速公路换电站运营由此成为一个具有用户车辆电池状态和多站点电池状态时空耦合、动态演化特征的序贯决策问题。

实际运营中存在两类异质用户：预约用户和随机到达用户。预约用户在运营日前提交计划到达高速入口的时间、行程起点和终点以及到达 SOC，运营者据此完成接纳决策并发布日前换电路径；随机用户直接到站提出请求，其到达时刻与 SOC 只能预测。预约用户的实际到达信息可能偏离预约内容；车辆进入高速公路后，其未执行的剩余换电站序列需要随实时位置、实时 SOC 和预计到达时间不断更新；到站时若暂无满电电池可用，用户在站内等待，超过其等待耐心后流失。考虑到分时电价使站内充电成本随时间变化，如何协调预约用户履约、随机用户接纳以及站内电池充电调度，是高速公路换电网络运营中的核心问题。

针对上述问题，本文构建 MPC--RL 分层协同优化框架。换电站的实时运营中存在两类性质不同但相互耦合的决策：一类是预约用户剩余换电路径的动态调整决策，即在有限预测域内根据未到用户的预约信息和在途用户的实时状态调整其未执行的换电站序列；另一类是各站逐槽电池的充电功率决策。剩余路径调整具有离散组合结构，并同时受到车辆续航、路径连通性、预约履约和多站点库存约束；随机需求预测与实际到达又随时间不断更新。模型预测控制（Model Predictive Control, MPC）能够在有限预测域内联合评估这些约束和未来库存影响，并在每次获得新观测后重新优化而只执行当前外层区间，因而适合处理该类动态路径调整问题。充电功率的影响则具有明显的跨时段累积性：当前功率不仅决定当期充电成本，还会改变未来各站逐槽电池的 SOC、满电电池供给能力和后续服务水平。本文采用强化学习（Reinforcement Learning, RL）学习逐槽充电功率控制策略，使其能够在长期交互过程中根据需求、电价和电池状态自适应调整充电行为。如\autoref{fig:MPC-RL}所示，RL actor 输出预测域内的逐槽请求充电功率并作为 MPC 的已知参数，MPC 在给定功率下联合优化剩余路径与站内连续服务事件安排，critic 向 MPC 提供预测域之外的长期价值信息。

\begin{figure}[!htbp]
  \centering
  % FIGURE_PROMPT_BEGIN: fig:MPC-RL
  % 绘制一张适用于运筹优化与智能交通学术论文的 MPC–RL 协同框架图。白色背景，扁平化矢量风格，横向 16:9，配色仅使用深蓝、青绿色、橙色和中性灰，无渐变、阴影、3D 或装饰图标。
  % 图中从左到右展示：当前观测状态（逐槽电池 SOC、预约与随机等待队列及剩余耐心、未到和在途预约信息、实时位置与 ETA、电价和随机需求预测）；基于已知信息预先确定的参考路径与参考执行轨迹；RL actor 输出预测域内逐槽请求功率，critic 输出终端 SOC 边际价值和边界未决交付价值；MPC 在给定 RL 信息后优化预约用户剩余路径和站内连续事件安排；确定性执行器在当前外层区间内按真实到站、充电完成、换电和超时事件连续执行；输出实际 reward、更新后的电池状态、等待队列和已发布剩余路径，并反馈至下一轮。
  % 明确区分外层固定间隔决策与区间内部连续事件，强调参考路径先确定、随后 actor rollout、最后 MPC 求解，不出现 RL–MPC 内循环。所有文字使用简洁中文，字体统一，箭头方向明确，模块不超过六个，适合缩放后在 A4 论文中阅读。输出优先为可编辑 SVG/PDF 或高分辨率透明背景图片。
  % FIGURE_PROMPT_END: fig:MPC-RL
  \fbox{\parbox[c][0.52\textwidth][c]{0.92\textwidth}{%
    \centering\small Figure placeholder}}
  \caption{MPC--RL 协同决策、连续事件执行与状态反馈框架}
  \label{fig:MPC-RL}
\end{figure}

进一步地，本文利用 RL critic 提供预测时域之外的长期价值信息，以缓解有限预测时域导致的末端效应。传统 MPC 通常通过设置固定终端库存下界避免时域末端过度消耗满电电池，但该方法只能粗略限制末端满电电池数量，难以区分各站逐槽电池 SOC 在未来需求和电价环境下的长期价值。本文在 RL 训练过程中学习状态价值函数，并在每轮 MPC 求解前计算 critic 对终端逐槽 SOC 的梯度，将其作为相应电池状态的长期边际价值。对于服务尚未发生、退回 SOC 尚未写入终端库存的未来预约交付，本文不额外训练新的 critic 分量，而是将其转化为“用退回电池替换一块满电电池”的等效终端影响。由此，RL 学到的长期电池状态价值被转化为 MPC 目标函数中的线性终端价值项。

因此，本文提出的 MPC--RL 框架基于决策属性进行分层：MPC 负责处理强约束、组合结构明显且需要可解释性的预约用户剩余路径动态调整问题；RL 负责学习具有长期影响的逐槽充电功率策略，并通过 critic 向 MPC 提供终端 SOC 价值信息。二者通过逐槽 SOC 状态、逐槽请求充电功率、在途预约履约事件和终端价值项相互连接，从而在保证预约用户服务质量要求和运营约束可行性的同时，使有限时域 MPC 能够近似感知预测时域之外的长期影响。

% CONTRIBUTIONS_PLACEHOLDER_BEGIN
% 待作者补充本文贡献
% CONTRIBUTIONS_PLACEHOLDER_END

本文其余部分安排如下：第 2 节给出问题描述；第 3 节建立滚动时域下的 MPC 模型，包括基于 SOC 分档的离线候选网络、日前基准路径生成和在线优化模型；第 4 节建立充电控制的马尔可夫决策过程并介绍 RL 算法；第 5 节给出数值实验；第 6 节总结全文。附录 A--D 分别给出候选事件构造与有效到达链、等待与优先级约束、连续充电与槽位匹配、路径发布与边界记账的实现细节。全文将站点规定的服务就绪 SOC 归一化为 $1$；因此，“满电”表示达到该运营目标值，并不要求等于电芯的电化学上限。






\section{Problem Statement}

考虑由多个 O-D 对及沿线换电站构成的单向高速公路换电网络。对于任一 O-D 对 $p$，车辆从入口 $o^p$ 进入高速公路后沿下游方向顺序行驶且不能折返，最终从出口 $d^p$ 离开；一次长途行程可能需要在沿线换电站集合 $\mathcal S^p$ 中的多个站点换电。站点 $i$ 配备槽位集合 $\mathcal B_i$，槽位中的电池以归一化 SOC 刻画电量状态。每次换电从站点取走一块 SOC 为 $1$ 的满电电池，并退回一块带有剩余电量的电池，退回电池入槽充电后可再次服务。因此，某位用户的站点选择不仅决定其后续可达站点，还会改变相关站点未来的电池库存，使用户路径与多站点库存产生时空耦合。

系统服务预约用户和随机到达用户。预约用户在运营日前提交预约信息，包括 O-D 对、预计进入高速公路的时刻及进入时的 SOC，运营者据此决定是否接纳预约，并为已接纳用户制定和发布日前换电路径。随机用户在到达换电站时提出请求，运营者只掌握其预测到达信息。按运行状态，已接纳预约用户分为两类：尚未进入高速公路的未到用户，以及已进入高速公路的在途用户；已经迟于预约入口时刻但尚未实际进入的用户仍属于未到集合，不能仅按原预约时刻将其删除。未到用户的有效入口 SOC 记为 $\widetilde{\mathrm{soc}}_{A,\ell}^{p,k}$，其初值为预约入口 SOC，并在获得新的入口前观测后更新；用户实际进入后，系统改用其实时位置、实时 SOC 和到达下游节点的预计到达时间。到站时若暂无满电电池可用，用户在站内等待：任一请求 $r$ 自其到站时刻 $t_r^{\mathrm{arr}}$ 起最长等待 $\Delta t$，开始服务时刻 $t_r^{\mathrm{sw}}$ 须满足 $t_r^{\mathrm{arr}}\le t_r^{\mathrm{sw}}\le t_r^{\mathrm{ddl}}$，其中 $t_r^{\mathrm{ddl}}=t_r^{\mathrm{arr}}+\Delta t$ 为等待截止时刻；截止时刻仍未获得服务的预约记为一次预约失败，随机请求则流失。换电视为瞬时事件。

系统以固定长度 $\sigma$ 的外层区间组织运营。以运营日统一时间原点，记
\[
t_q=q\sigma,
\qquad
\mathcal I_q=[t_q,t_{q+1}),
\qquad
\mathcal T_\ell=\{\ell,\ldots,\ell+H-1\}.
\]
在每个区间起点 $t_\ell$，系统以未来 $H$ 个区间为预测域，即 $\mathcal T_\ell$ 对应连续时间范围 $[t_\ell,t_{\ell+H})$。预约与随机请求的到站时刻均为外生连续值，区间索引只用于请求归组、电价与站级充电能量限额等分段常数数据。服务顺序遵循两条规则：预约请求优先于随机请求；同类请求内部按原到站时刻升序，同一时刻按固定用户编号升序。系统不为尚未到站的请求扣留当前可用电池；等待请求跨滚动边界携带，保留其请求标识、原到站时刻与原截止时刻，不重新计时。

运营者需要协调预约用户的剩余换电站序列与各站逐槽电池的充电功率，从而在履行已接纳预约的前提下，权衡换电收益、购电成本、路径调整成本和长期库存价值。本文采用分层 MPC--RL 框架：RL 根据需求、电价和库存状态生成预测域内的逐槽请求充电功率，并向 MPC 提供终端价值系数；MPC 在给定请求功率下联合优化剩余路径与站内连续服务事件安排。每轮优化后，系统发布发生变化的在途用户剩余站序，只执行首个外层区间内的充电与服务安排，随后根据真实到站、充电完成、换电和超时事件更新状态并进入下一轮求解。

\autoref{tab:notation}列出正文反复使用的核心符号；仅在附录使用的符号在各附录中就近说明。RL 的状态、动作和算法符号在第 4 节定义。

\begin{table}[!htbp]
  \centering
  \caption{正文主要符号}
  \label{tab:notation}
  \small
  \setlength{\tabcolsep}{3pt}
  \renewcommand{\arraystretch}{0.98}
  \begin{tabular}{@{}lp{11cm}@{}}
    \toprule
    \textbf{符号} & \textbf{含义} \\
    \midrule
    \multicolumn{2}{l}{\textit{集合与索引}} \\
    $\sigma,t_q,\mathcal I_q,\mathcal T_\ell,\Delta t$ & 外层区间长度与起点、半开区间、含 $H$ 个区间的预测域索引集，以及最大等待时间 \\
    $\mathcal P,p;\ \mathcal S^p,\mathcal S;\ \mathcal B_i,b$ & O-D 对、沿线与全系统站点、库存槽位的集合和索引；槽位不是永久电池编号 \\
    $\mathcal K_\ell^{\mathrm{fut},p},\mathcal K_\ell^{\mathrm{enr},p},\mathcal K_\ell^p$ & 未到、在途预约用户集合及其并集；迟到但尚未实际进入者仍属于未到集合 \\
    $o^p,d^p,o_\ell^{\mathrm{cur},p,k},\widetilde{\mathrm{soc}}_{A,\ell}^{p,k},\mathrm{soc}_{A,\ell}^{p,k}$ & 入口、出口、实时虚拟起点、未到用户有效入口 SOC 与在途用户实时 SOC \\
    $\mathcal A^{p,k},t_{A,i}^{p,k},t_r^{\mathrm{arr}},t_r^{\mathrm{ddl}},t_r^{\mathrm{sw}}$ & 个体剩余候选弧集、连续 ETA，以及请求的到站、截止和开始服务时刻 \\
    \midrule
    \multicolumn{2}{l}{\textit{参数与状态}} \\
    $D_{i,j}^{p},v_{i,j}^{p},\underline D,\underline s_d^p$ & 节点间距离与 SOC 消耗、最小换电间距和出口最低 SOC \\
    $e_{i,q},\pi_{i,q}^{\mathrm{sw}},\overline E_{i,q},\overline P_i$ & 购电价、服务价格、站级区间充电能量上限和单槽功率上限 \\
    $E_B,\eta_i,\widehat P_{i,b,q},P_{i,b}(t),S_{i,b}(t)$ & 电池额定能量、充电效率、RL 请求功率、实际功率和槽位 SOC \\
    $\bar{\mathbf y}_{A}^{p,k},\mathbf y_{A}^{\mathrm{pub},p,k}$ & 日前初始发布路径与最近一次实际发布的剩余路径 \\
    $\kappa,c^{\mathrm{fail}},\beta$ & 单次路径调整成本、单次预约失败成本与 RL 终端价值权重 \\
    \midrule
    \multicolumn{2}{l}{\textit{决策变量与主要辅助量}} \\
    $y_{i,j}^{p,k},x_{A,i}^{p,k},d_{p,k}$ & 剩余路径、站点访问量与路径差异指示量 \\
    $s_\varepsilon^A,f_\varepsilon^A,\omega_\varepsilon^A,z_r,\omega_r^R$ & 预约和随机请求的服务、超时或流失、跨边界未决结果 \\
    $\lambda_{i,b}^S,\Delta V_i^{\mathrm{out}}(\rho),w_h^{\mathrm{pend}}$ & 终端 SOC 边际价值、未实现交付的等效价值及其未决指示量 \\
    \bottomrule
  \end{tabular}
\end{table}


\section{MPC Formulation}

设当前决策区间为 $\mathcal I_\ell$。在时刻 $t_\ell$，系统读取：各站逐槽电池的观测 SOC 与正在充电的状态；自上一轮携带而来、仍未超时的站内等待请求（保留其原请求标识、原到站时刻与原截止时刻）；在途预约用户的实时位置、实时 SOC 与到达下游节点的预计到达时间；已发布且尚未执行的剩余路径；以及未来 $H$ 个区间的随机需求预测与电价。

随后，系统按固定顺序构造本轮求解输入。首先确定参考路径：在途用户取最近一次实际发布路径的未执行站序；未到用户优先沿用上一轮保留的内部参考路径，该路径不存在或在当前信息下不再满足几何与续航可行性时，以日前初始发布路径为基准修复。参考路径修复只使用求解前已知的位置、SOC、预计到达时间和候选网络，不依赖尚未生成的 actor 功率或本轮优化变量。其次，以同一 actor 的均值策略沿参考执行轨迹滚动生成预测域内的逐槽请求充电功率 $\widehat P_{i,b,q}$，并由 critic 计算终端价值系数（第 4 节）。参考执行采用与实际执行相同的优先级、等待、超时、充电与槽位替换规则，并允许预约软失败。因此，请求功率序列与终端价值系数同本轮待求的路径、服务和匹配变量无关，不构成循环依赖。

在此基础上，系统为每位预约用户构造本轮剩余候选网络（第 3.1 节），并由 MPC 在给定请求功率下联合优化剩余路径与站内连续事件安排（第 3.3 节）。模型采用半开预测域 $[t_\ell,t_{\ell+H})$：恰好发生在 $t_{\ell+H}$ 的到站、充满、服务或超时事件由下一轮处理；终端 SOC 取事件前快照 $S_{i,b}(t_{\ell+H}^{-})$，端点事件不跨轮重复执行或结算。

求解后，系统立即发布发生变化的在途用户剩余站序并结算一次路径调整成本，新站序作为下一轮的调整基准；未到用户的内部优化路径不发布、不覆盖日前路径，也不进入当前实际收益。系统只执行首个外层区间 $\mathcal I_\ell$：区间内按真实到站、充电完成、换电和超时事件连续推进；区间结束时，逐槽 SOC、仍未超时的等待队列与已发布剩余路径进入下一轮。这一机制可概括为“滚动时域优化、连续事件执行”，如\autoref{fig:rolling_horizon}所示。

\begin{figure}[!htbp]
  \centering
  % FIGURE_PROMPT_BEGIN: fig:rolling_horizon
  % 绘制一张适用于学术论文的横向滚动时域示意图。白色背景、简洁矢量风格，宽高比约 2:1，使用深蓝表示预测域、橙色表示当前执行区间、绿色表示跨边界携带状态、灰色表示已执行历史。
  % 顶部为时间轴，清晰标出 $t_{\ell-1}$、$t_\ell$、$t_{\ell+1}$ 和 $t_{\ell+H}$，预测域采用半开区间 $[t_\ell,t_{\ell+H})$，突出只执行首个外层区间。下方设置三条泳道：第一条为未到预约用户，显示其内部路径可滚动优化但仍相对初始发布路径评价；第二条为在途预约用户，显示从实时虚拟起点重优化剩余站序，路径变化后发布并更新比较基准；第三条为站内等待队列，显示原始到站时刻和截止时刻跨滚动边界保持不变。
  % 当前执行区间内部使用少量时间标记展示连续到站、充电完成、立即换电和超时事件；预测终点处显示尚未截止的请求作为未决状态传入下一轮。不得出现“已发布路径固定不再优化”“整时段统一换电”或“随机用户必须当期完成”等旧逻辑。文字精炼、层级清楚、无大段说明，适合 A4 论文单栏宽度阅读。输出优先为可编辑 SVG/PDF 或高分辨率图片。
  % FIGURE_PROMPT_END: fig:rolling_horizon
  \fbox{\parbox[c][0.46\textwidth][c]{0.92\textwidth}{%
    \centering\small Figure placeholder}}
  \caption{滚动时域中的剩余路径优化、连续事件执行与等待请求携带}
  \label{fig:rolling_horizon}
\end{figure}


\subsection{Offline Candidate networks by SOC Range}

在每个决策时刻，需要安排预测域内每位预约用户的剩余换电站序列。如何高效、清晰地表示解的结构并设计决策变量是首要问题。本文借鉴加油站选址模型中 expanded network 的思路，用一条入口--出口路径表示该序列：对于 O-D 对 $p$，若路径为
$o^p\rightarrow i_1\rightarrow\cdots\rightarrow i_r\rightarrow d^p$，则用户依次在换电站 $i_1,\ldots,i_r$ 换电，一条弧跨过的沿线站点均不被访问。为避免在每个滚动时刻重复枚举全部下游组合，本文在运营前先生成候选路径网络。不同的用户出发 SOC 会影响车辆的续航里程，从而影响可达的换电站集合。我们将用户出发 SOC 划分为若干个离线区间，并在每个区间上生成相应的候选网络。

对每个 O-D 对 $p$，节点集合为
$\mathcal V^p=\{o^p\}\cup\mathcal S^p\cup\{d^p\}$，所有节点均按高速公路下游方向排列。参数 $D_{i,j}^p$ 和 $v_{i,j}^p>0$ 分别表示从 $i$ 到下游节点 $j$ 的行驶距离和 SOC 消耗。

考虑到现实场景中为用户安排两个距离过近的连续站点换电并不合理，同时也为了减少网络规模，给定最小换电间距 $\underline D>0$，生成的网络在保证方案可行性的基础上，尽量排除从出发位置到首个换电站以及任意两个连续访问换电站之间距离小于 $\underline D$ 的方案。指向出口的末段行程不受该间距约束，但车辆到达出口时的 SOC 不得低于 $\underline s_d^p$。

对每个 SOC 区间 $h$ 和 O-D 对 $p$，离线网络按以下步骤生成。

\begin{enumerate}
  \item[\textbf{Step 1.}] \textbf{生成原始弧。}若车辆以该 SOC 区间的上限电量能够从入口 $o^p$ 到达下游换电站 $j$，则连接二者；若 $v_{i,j}^p\le1$，则连接两个依次位于下游的换电站 $i$ 和 $j$；若 $v_{i,d^p}^p+\underline s_d^p\le1$，则连接换电站 $i$ 与出口 $d^p$。

  \item[\textbf{Step 2.}] \textbf{生成完整可行路径。}在原始网络中枚举所有从 $o^p$ 到 $d^p$ 的完整路径。

  \item[\textbf{Step 3.}] \textbf{剪枝过近换电弧。}从上述完整路径中找出距离小于 $\underline D$ 的首段弧和站间弧，并按距离从小到大排序。依次考察每条过近弧，如果现有的完整路径中仍存在至少一条不包含该弧的方案，则删除该弧，否则保留。

  \item[\textbf{Step 4.}] \textbf{形成候选网络。}将最终保留的完整路径所包含的弧合并为 $\mathcal A^{h,p}$，候选网络由节点集 $\mathcal V^p$ 和弧集 $\mathcal A^{h,p}$ 组成。
\end{enumerate}

离线分档网络提供下游站间连接骨架。在线使用时，未到用户以其有效入口 SOC $\widetilde{\mathrm{soc}}_{A,\ell}^{p,k}$ 所在区间对应的离线网络为基础；对在途用户，系统删除其已经经过的节点和弧，并从实时虚拟起点 $o_\ell^{\mathrm{cur},p,k}$ 向下游节点补建首段弧。虚拟首段依据求解前已知的实时位置、实时 SOC 和连续预计到达时间筛选可达性，本文不为此建立随机交通模型。最小换电间距 $\underline D$ 始终以入口或最近一次真实换电位置为基准解释，不因为虚拟起点随滚动前移而把用户视为刚刚完成一次换电。由此得到的个体剩余候选弧集记为 $\mathcal A^{p,k}$；已执行的路径前缀不进入本轮决策。


\subsection{Day-Ahead Baseline Path Generation}

用户提交预约信息
$\left(p,\bar t_A^{p,k},\overline{\mathrm{soc}}_A^{p,k}\right)$
后，运营者依据其行程、预计入口时刻和入口 SOC 检查能否提供服务，并为可接纳的请求生成日前基准路径。

预约请求按提交顺序处理。处理当前请求前，算法根据已接纳预约和日前随机需求预测，递推估计各站未来各时段的槽位库存。日前充电功率由训练完成的同一 actor 的均值策略在日前预测状态上逐区间生成；库存递推采用与在线模型一致的连续充电规则：电池未满时按当前区间的请求功率充电，达到 SOC $1$ 后立即停止充电，换电退回的低 SOC 电池入槽后立即恢复充电；预计换电事件按预计服务时刻扣减满电库存。因此，给定初始库存、需求预测和下述并列规则，日前库存递推与预约接纳结果可以复现。

对 O-D 对 $p$ 的预约用户 $k$，基准路径按以下规则生成：
\begin{enumerate}
  \item 车辆从入口出发时的 SOC 为 $\overline{\mathrm{soc}}_A^{p,k}$，每次换电后恢复为 $1$。下一节点沿候选弧集选择，因此首段和站间行程均满足续航要求。
  \item 若车辆能够从当前位置直接到达出口，且到达 SOC 不低于 $\underline s_d^p$，则下一站选择出口；否则，仅保留预计到站时刻已经有满电电池可用、并且仍存在可行下游延续路径的候选站。该接纳标准是严格的：预计到达时必须已经存在满电电池，不利用等待时间 $\Delta t$ 扩大接纳量。
  \item 候选站按预计满电电池余量从高到低排序，余量相同时优先选择更靠近出口的站点。选定站点后，在相应时段为该用户预留一块满电电池，并更新该站后续时段的预计库存；随后重复上述过程。
\end{enumerate}

若任一步不存在满足库存与下游可达性要求的候选站，即不能生成完整入口--出口路径，则拒绝该请求；只有成功生成完整路径的用户才进入已接纳集合。对已接纳用户，若基准路径经过弧 $(i,j)$，则令 $\bar y_{A,i,j}^{p,k}=1$，否则为 $0$。该路径随预约结果一并发布，既作为用户的初始发布路径，也作为其未到阶段路径调整的固定比较基准。


\subsection{Optimization Model}

给定参考路径、逐槽请求充电功率、终端价值系数与个体剩余候选网络，MPC 在预测域内联合优化预约用户的剩余路径与站内连续服务事件安排。

\paragraph{Decision variables.}
路径决策变量为
\[
y_{i,j}^{p,k}\in\{0,1\},
\qquad
(i,j)\in\mathcal A^{p,k},
\quad
k\in\mathcal K_\ell^p.
\]
若 O-D 对 $p$ 的预约用户 $k$ 经过弧 $(i,j)$，则 $y_{i,j}^{p,k}=1$；一组满足流平衡的 $y$ 变量唯一确定该用户剩余行程中依次访问的换电站序列。路径变量属于本轮子问题，不随滚动轮次增加下标。辅助变量 $x_{A,i}^{p,k}$ 和 $d_{p,k}$ 分别表示站点访问量和路径差异指示量。

站内服务以通用请求索引 $r$ 描述。每个请求携带连续的开始服务时刻 $t_r^{\mathrm{sw}}$；对预约候选事件 $\varepsilon$（定义见附录 A），二元变量 $s_\varepsilon^A$、$f_\varepsilon^A$ 和 $\omega_\varepsilon^A$ 分别表示其在预测域内获得服务、等待超时和跨边界未决；对随机请求 $r$，二元变量 $z_r$ 和 $\omega_r^R$ 分别表示其获得服务和跨边界未决，随机流失由 $1-z_r-\omega_r^R$ 派生，不单独设变量。槽位 SOC 轨迹 $S_{i,b}(t)$ 为连续时间状态量。

逐槽请求充电功率 $\widehat P_{i,b,q}$ 是 RL actor 经参考 rollout 生成的已知参数（第 4 节）。实际充电功率 $P_{i,b}(t)$ 不是 MPC 自由选择的变量，而由确定性充电规则派生：电池未满时等于当前区间的请求功率，达到 SOC $1$ 后立即停止充电；同一槽位换入低 SOC 电池后，立即恢复当前区间的同一请求功率。附录 B--C 以站级物理时刻 $m$ 展开队列与槽位状态；$x_{i,r,m}$ 和 $\alpha_{i,b,r,m}$ 分别表示请求在该时刻获得服务及其槽位匹配，事件顺序与充满时刻由确定性递推生成而不能被优化器任意延迟。

\paragraph{Objective function.}
本轮 MPC 的目标按服务收益、连续充电成本、路径调整成本、预约超时成本和 RL 终端价值依次定义。

\textbf{换电服务收益。}令 $\pi_{i,q}^{\mathrm{sw}}\ge0$ 表示站点 $i$ 在区间 $q$ 的单位换电电量服务价格。用户获得满电电池并退回 SOC 为 $\rho$ 的电池时，本次换电的服务电量为 $E_B(1-\rho)$。服务收益按预测服务时刻所在价格区间取值：
\begin{equation}
I^{\mathrm{MPC}}
=E_B\sum_{i\in\mathcal S}
\sum_{r\in\mathcal E_i}
\sum_{q\in\mathcal T_\ell}
\pi_{i,q}^{\mathrm{sw}}(1-\rho_r)\gamma_{i,r,q},
\label{eq:income}
\end{equation}
其中 $\mathcal E_i$ 为站点 $i$ 的预约、随机及携带请求记录集合，$\rho_r$ 为其退回 SOC；$\gamma_{i,r,q}=1$ 当且仅当请求 $r$ 在半开区间 $\mathcal I_q$ 内获得服务，其事件位置线性化见附录 B。将请求集合按类别拆分即可得到 $I^{\mathrm{MPC}}=I^A+I^R$；$I^A$ 只计入由站内满电库存实际履约的预约换电，$I^R$ 使用当前预测到站信息，因而是前瞻收益。个体候选网络已规定车辆一旦能够满足出口 SOC 要求便直接驶向出口，因此正的服务收益不会诱导可避免的额外换电。

\textbf{连续充电成本。}预测域内的购电成本按实际功率积分计算：
\begin{equation}
C_{\mathrm{ch}}
=
\sum_{i\in\mathcal S}
\sum_{b\in\mathcal B_i}
\int_{t_\ell}^{t_{\ell+H}}
e_i(t)P_{i,b}(t)\,\mathrm dt,
\label{eq:chargingcost}
\end{equation}
其中 $e_i(t)=e_{i,q}$（$t\in\mathcal I_q$）为分段常数电价。给定请求功率后，相邻事件之间的充电能量是“已知功率系数乘以连续持续时间”，不产生两个自由连续变量的乘积；事件级展开见附录 C。

\textbf{路径调整成本。}若预约用户的剩余换电站序列不同于其比较基准，则在发布时产生一次调整成本：
\begin{equation}
C_{\mathrm{adj}}
=
\kappa
\sum_{p\in\mathcal P}
\sum_{k\in\mathcal K_\ell^p}
d_{p,k}.
\label{eq:adjustment_cost}
\end{equation}
每名用户无论改变一个还是多个换电站均只计一次惩罚。未到用户的比较基准是日前初始发布路径，在途用户的比较基准是最近一次实际发布路径的未执行站序；虚拟起点随滚动前移而剩余站序不变时不产生调整成本。未到用户的调整成本是前瞻评价；只有本轮真实发布调整的在途用户产生实际支出。

\textbf{预约超时成本。}当有效预约事件在其等待截止时刻仍未获得服务时，记一次预约失败：
\begin{equation}
C_{\mathrm{fail}}
=
c^{\mathrm{fail}}
\sum_{i\in\mathcal S}
\sum_{\varepsilon\in\mathcal E_{i}^{A}}
f_\varepsilon^A.
\label{eq:reservation_failure_cost}
\end{equation}
该项以预约换电事件为结算单位，失败成本按截止时刻所在的执行区间结算。预测域内的 $f_\varepsilon^A$ 表示当前信息下的超时风险并进入本轮目标；只有真实发生的超时才进入 RL reward 和全天实际汇总，重叠滚动窗口中的预测风险不重复结算。模型不设置等待时间成本，也不设置随机流失成本。

\textbf{RL 终端价值。}有限预测域会截断两类长期影响：域内决策在时刻 $t_{\ell+H}$ 留下的槽位 SOC 会继续影响后续运营；本轮所选路径中尚未在预测边界前完成服务的预约交付也会在未来消耗满电库存。后一类既包括预计在边界之后到达的交付，也包括已经到站但在边界处仍未超时、仍在等待的请求。为评价这两类影响，本文使用第 4 节训练的 critic $V_\theta$ 为 MPC 提供预计算的价值系数，而不将非线性网络直接嵌入优化模型。

每轮求解前，系统依据参考路径、当前预约信息、随机需求预测和 actor 的均值动作生成一条与本轮路径变量无关的参考执行轨迹，其终端状态记为 $s_{\ell+H}^{\mathrm{ref}}$，包含逐槽 SOC 与仍未超时的等待队列。在该状态处，critic 对槽位 SOC 的一阶边际价值为
\begin{equation}
\lambda^S_{i,b}
=
\left.
\frac{\partial V_{\theta}(s)}
{\partial S_{i,b}}
\right|_{s=s_{\ell+H}^{\mathrm{ref}}},
\quad i\in\mathcal S,\ b\in\mathcal B_i.
\label{eq:marginal_value_soc}
\end{equation}
对 critic 作一阶近似并删除与本轮决策无关的常数后，域内终端库存只需保留线性项 $\lambda^S_{i,b}S_{i,b}(t_{\ell+H}^{-})$。

对于任一尚未在预测边界前完成服务、因而其退回 SOC 尚未写入终端库存的候选预约交付，站点 $i$ 将在后续运营中发出一块满电电池并接收 SOC 为 $\rho$ 的退回电池。该交付需要与域内终端 SOC 分开估值。为使反事实在参考终端库存暂时没有满电电池时仍有定义，在站点 $i$ 选择一个规范化槽位：若存在服务就绪槽位，则取编号最小者；否则取 SOC 最高者，并以最小编号打破并列。记 $s_{\ell+H}^{\mathrm{ref},(i,x)}$ 为仅将该槽位 SOC 置为 $x$、其余状态保持不变的参考终端状态，并定义
\begin{equation}
\Delta V_{i}^{\mathrm{out}}(\rho)
=
 V_\theta\!\left(s_{\ell+H}^{\mathrm{ref},(i,\rho)}\right)
-V_\theta\!\left(s_{\ell+H}^{\mathrm{ref},(i,1)}\right).
\label{eq:outside_swap_value}
\end{equation}
$\Delta V_{i}^{\mathrm{out}}(\rho)$ 在 MPC 求解前由 critic 计算，作为固定系数。当 $\rho<1$ 且 critic 合理反映库存价值时，该值通常为负，对应未来履约的库存机会成本。上述成对反事实比较的是同一规范化槽位从满电变为退回 SOC 的价值变化，不要求参考轨迹事先在该站保留满电电池。

去除不影响最优解的常数项后，进入综合目标的 RL 终端价值为
\begin{equation}
\Phi^{\mathrm{RL}}
=
\sum_{i\in\mathcal S}\sum_{b\in\mathcal B_i}
\lambda^S_{i,b}S_{i,b}(t_{\ell+H}^{-})
+\sum_{h\in\overline{\mathcal D}_{\ell}^{A}}
\Delta V_{i(h)}^{\mathrm{out}}(\rho_h)
w_h^{\mathrm{pend}},
\label{eq:rl_terminal_value}
\end{equation}
其中 $\overline{\mathcal D}_{\ell}^{A}$ 是由本轮候选弧和携带预约记录预先枚举的候选交付记录集，$w_h^{\mathrm{pend}}$ 表示交付 $h$ 在预测边界仍未物理实现：相应路径被选择或携带请求有效、用户未因上游超时而失活，且该交付尚未在域内完成时取 $1$。已在域内服务并写入终端 SOC 的交付对应指示为 $0$，每个未来交付只估值一次。边界存活与防重复估值的完整展开见附录 D。多个未决交付的价值采用可加近似，忽略其高阶交互；$\beta$ 用于校准 critic 价值与域内运营金额的尺度。

综合上述各项，本轮 MPC 的目标为最大化综合运营价值：
\begin{equation}
\max\quad
I^{\mathrm{MPC}}-C_{\mathrm{ch}}-C_{\mathrm{adj}}-C_{\mathrm{fail}}
+\beta\Phi^{\mathrm{RL}}.
\label{eq:objective}
\end{equation}

\paragraph{Constraints.}
正文给出六组代表性关系；候选事件枚举、有效到达链、等待与优先级、槽位匹配和边界条件的实现级展开分别在附录 A--D 给出。

（1）\textbf{剩余路径流平衡与续航可达。}每位预约用户满足流平衡约束：
\begin{equation}
\sum_{\{j\mid(n,j)\in\mathcal A^{p,k}\}}
y_{n,j}^{p,k}
-
\sum_{\{j\mid(j,n)\in\mathcal A^{p,k}\}}
y_{j,n}^{p,k}
=
\begin{cases}
1,&n=o_\ell^{p,k},\\
-1,&n=d^p,\\
0,&\text{其余节点},
\end{cases}
\quad
p\in\mathcal P,\ k\in\mathcal K_\ell^p,\ n\in\mathcal V^{p,k},
\label{eq:flow}
\end{equation}
其中本轮起点 $o_\ell^{p,k}$ 对未到用户为入口 $o^p$，对在途用户为实时虚拟起点 $o_\ell^{\mathrm{cur},p,k}$；$\mathcal V^{p,k}$ 为该用户剩余候选网络的节点集。续航可达性由候选网络构造保证：首段弧按有效入口 SOC（未到用户）或实时 SOC（在途用户）筛选，站间弧以满电出发，出口弧满足 $\underline s_d^p$。用户 $k$ 沿入弧 $(j,i)$ 到达站点 $i$ 时，退回电池的 SOC 为
\begin{equation}
\rho_{A,j,i}^{p,k}
=
\begin{cases}
\widetilde{\mathrm{soc}}_{A,\ell}^{p,k}-v_{o^p,i}^{p},
&j=o^p,\ k\in\mathcal K_\ell^{\mathrm{fut},p},\\[2pt]
\mathrm{soc}_{A,\ell}^{p,k}
-v_{o_\ell^{\mathrm{cur},p,k},i}^{p,k},
&j=o_\ell^{\mathrm{cur},p,k},\ k\in\mathcal K_\ell^{\mathrm{enr},p},\\[2pt]
1-v_{j,i}^{p},
&j\in\mathcal S^p,
\end{cases}
\quad (j,i)\in\mathcal A^{p,k}.
\label{eq:return_soc}
\end{equation}
即未到用户首段使用当前有效入口 SOC，在途用户首段使用实时 SOC，后续站间弧使用上一站成功换得的满电 SOC。站内等待中的在途用户保留其当前请求与当前站点，路径重优化不能取消该等待请求。

（2）\textbf{按用户状态区分基准的路径调整识别。}定义当前路径的站点访问量
\begin{equation}
x_{A,i}^{p,k}
=
\sum_{\{j\mid(j,i)\in\mathcal A^{p,k}\}}y_{j,i}^{p,k},
\qquad i\in\mathcal S^p,
\label{eq:station_visit_indicator}
\end{equation}
以及参考访问量
\begin{equation}
x_{A,i\mid\ell}^{\mathrm{ref},p,k}
=
\begin{cases}
\bar x_{A,i}^{p,k},
&k\in\mathcal K_\ell^{\mathrm{fut},p},\\
x_{A,i\mid\ell}^{\mathrm{pub},p,k},
&k\in\mathcal K_\ell^{\mathrm{enr},p},
\end{cases}
\label{eq:reference_visit_indicator}
\end{equation}
其中 $\bar x_{A,i}^{p,k}$ 由日前初始发布路径 $\bar y_A$ 汇总，$x_{A,i\mid\ell}^{\mathrm{pub},p,k}$ 由最近一次实际发布路径的未执行站序汇总（附录 D）。比较范围为当前候选实体站点与参考未执行站点的并集，缺失访问量补零；已执行的历史前缀不参与比较。路径差异指示量满足
\begin{subequations}
\label{eq:path_adjustment_indicator}
\begin{align}
d_{p,k}
&\ge x_{A,i}^{p,k}-x_{A,i\mid\ell}^{\mathrm{ref},p,k},
&&p\in\mathcal P,\ k\in\mathcal K_\ell^p,\ i\in\mathcal S^p,\\
d_{p,k}
&\ge x_{A,i\mid\ell}^{\mathrm{ref},p,k}-x_{A,i}^{p,k},
&&p\in\mathcal P,\ k\in\mathcal K_\ell^p,\ i\in\mathcal S^p.
\end{align}
\end{subequations}
只要剩余站序相对其基准增加或删除任意换电站，就有 $d_{p,k}=1$；与基准一致时，在目标函数中的正调整成本作用下有 $d_{p,k}=0$。

（3）\textbf{请求服务窗与服务/超时/未决守恒。}通用请求 $r$ 的等待截止时刻为
\begin{equation}
t_r^{\mathrm{ddl}}
=
t_r^{\mathrm{arr}}+\Delta t,
\label{eq:waiting_deadline}
\end{equation}
获得服务的请求满足 $t_r^{\mathrm{arr}}\le t_r^{\mathrm{sw}}\le t_r^{\mathrm{ddl}}$。预约候选事件满足守恒关系
\begin{equation}
s_\varepsilon^A+f_\varepsilon^A+\omega_\varepsilon^A=a_\varepsilon^A,
\label{eq:reservation_conservation_main}
\end{equation}
其中 $a_\varepsilon^A$ 为考虑上游服务链后的有效到达量（附录 A）：上游预约超时后，下游预约事件失活，只结算真实发生的一次失败。对随机请求令 $h_r^T=\mathbb I\{t_r^{\mathrm{ddl}}\ge t_{\ell+H}\}$，并施加 $\omega_r^R\le h_r^T$、$z_r+\omega_r^R\le1$ 和 $z_r+\omega_r^R\ge h_r^T$；因此，等待窗跨预测终点时，请求只能在域内服务或保持未决，不能提前流失，完整等待窗位于域内时才由 $1-z_r-\omega_r^R$ 派生流失。跨边界携带的请求保留原请求标识、原到站时刻与原截止时刻。服务窗、守恒与边界条件的完整约束见附录 B。

（4）\textbf{连续 SOC 积分与换电跳变。}在相邻槽位事件时刻 $t_1<t_2$ 之间，槽位 SOC 满足
\begin{equation}
S_{i,b}(t_2)-S_{i,b}(t_1)
=
\frac{\eta_i}{E_B}
\int_{t_1}^{t_2}P_{i,b}(t)\,\mathrm dt,
\qquad
P_{i,b}(t)=
\begin{cases}
\widehat P_{i,b,q(t)},&S_{i,b}(t)<1,\\
0,&S_{i,b}(t)=1,
\end{cases}
\label{eq:event_charging_transition}
\end{equation}
即电池未满时按当前区间请求功率充电，达到 SOC $1$ 后立即停止。请求功率的可行投影为
\begin{equation}
0\le \widehat P_{i,b,q}\le\overline P_i,
\qquad
\sigma\sum_{b\in\mathcal B_i}\widehat P_{i,b,q}
\le\overline E_{i,q},
\qquad i\in\mathcal S,\ q\in\mathcal T_\ell.
\label{eq:power_projection}
\end{equation}
服务就绪阈值严格为 SOC $1$；求解器数值容差不进入能量、收入或续航公式。换电为瞬时事件，服务完成后槽位 SOC 立即写入退回 SOC：$S_{i,b}\big((t_r^{\mathrm{sw}})^+\big)=\rho_r$；槽位 $b$ 是库存槽位而非永久电池编号，同一槽位在一个外层区间内可多次服务，但任意两次交付之间必须重新达到 SOC $1$。包含全部区间边界的站级物理时刻递推、严格满电标记和线性槽位匹配见附录 C。

（5）\textbf{预约优先、类内顺序与服务就绪。}在每个服务时刻，只对已到站、仍在等待且未超时的请求分配满电电池：预约请求优先于随机请求，同类请求内部按原到站时刻升序、同一时刻按固定编号升序；每次服务使用编号最小的服务就绪槽位；不为尚未到站的预约请求扣留当前可用电池。同一时刻的固定处理顺序为：充电完成、登记到站、队列分配、对仍未服务者判超时；截止时刻恰好有电池充满时，先允许服务再判定超时。附录 B 以队列流、优先前缀和功守恒约束强制上述服务纪律，附录 C 展开并行槽位容量与最小就绪槽位匹配。

（6）\textbf{终端未决、初始状态与变量范围。}本轮初始槽位 SOC、正在充电的状态和携带队列均取真实观测；初始条件、SOC 范围和槽位事件域在附录 C 给出，请求结果与半开服务域在附录 B 给出。终端状态取事件前快照 $S_{i,b}(t_{\ell+H}^{-})$ 与仍未超时的等待队列，未决交付及其防重复估值约束在附录 D 给出。模型不设置预测域末端的满电电池数量下界，域外长期影响统一通过式\eqref{eq:rl_terminal_value}进入综合目标。



\section{RL Formulation}

RL 层用于学习逐槽充电功率控制策略，并通过 critic 为 MPC 提供终端逐槽 SOC 价值信息。本节将充电功率控制问题建模为马尔可夫决策过程，并给出本文采用的 RL 算法。

\subsection{Markov decision process}

\paragraph{State.}
在滚动决策时刻 $t_\ell$，RL agent 观测系统状态 $s_\ell$。除各站逐槽电池的 SOC、电价和时间特征外，状态包含：未到与在途预约用户信息（在途用户含实时位置、实时 SOC 和最近发布的剩余站序）、站内预约与随机等待队列（含原到站时刻与剩余耐心）、最近一个外层区间的真实随机到达记录，以及未来逐请求随机需求预测。令
$\mathcal K_\ell^{\mathrm{rem},p}$ 表示时刻 $t_\ell$ 尚未进入高速公路的已接纳预约用户集合（覆盖范围超出当前预测域），用 $\bar{\mathbf y}_A^{p,k}$ 和 $\mathbf y_A^{\mathrm{pub},p,k}$ 分别汇集日前初始发布路径与最近实际发布路径的弧指示量，则预约信息集合为
\begin{equation}
\begin{aligned}
\mathcal Z_{A}
={}&
\left\{
\left(p,k,\bar t_A^{p,k},
\overline{\mathrm{soc}}_A^{p,k},
\bar{\mathbf y}_A^{p,k}\right):
k\in\mathcal K_\ell^{\mathrm{rem},p}
\right\}\\
&\cup
\left\{
\left(p,k,o_\ell^{\mathrm{cur},p,k},\mathrm{soc}_{A,\ell}^{p,k},
\mathbf y_A^{\mathrm{pub},p,k}\right):
k\in\mathcal K_\ell^{\mathrm{enr},p}
\right\}.
\end{aligned}
\label{eq:rl_reservation_information}
\end{equation}
站内等待队列记为
$\mathcal W_\ell=\{(r,i,t_r^{\mathrm{arr}},t_r^{\mathrm{ddl}})\}$，
其中元素为仍未服务且未超时的预约与随机请求。由于 $\mathcal Z_{A}$ 与 $\mathcal W_\ell$ 的元素数量随滚动时刻变化，分别使用 $f_A$、$f_W$ 将其编码为固定维度特征。再令 $f_B$、$f_R^{\mathrm{act}}$ 和 $f_R^{\mathrm{pred}}$ 分别编码逐槽 SOC、真实随机到达记录和预测随机请求，则状态表示为
\begin{equation}
\begin{aligned}
s_\ell=\Big(&
f_B\!\left(\{S_{i,b}^{\mathrm{obs}}(t_\ell)\}_{i\in\mathcal S,b\in\mathcal B_i}\right),
f_A(\mathcal Z_{A}),
f_W(\mathcal W_\ell),\\
&
f_R^{\mathrm{act}}\!\left(
\{(i,r,t_R^{i,r},\mathrm{soc}_R^{i,r},z_{R,r,i}^{\mathrm{act}}):
r\in\mathcal R_{i,\ell-1}^{\mathrm{act}}\}_{i\in\mathcal S}
\right),\\
&
f_R^{\mathrm{pred}}\!\left(
\{(i,\nu,t_{R}^{i,\nu},\mathrm{soc}_{R}^{i,\nu}):
\nu\in\mathcal R_{i,q}\}_{i\in\mathcal S,q\in\mathcal T_\ell}
\right),
\{e_{i,q}\}_{i\in\mathcal S,q\in\mathcal T_\ell},
\xi_\ell
\Big).
\end{aligned}
\label{eq:rl_state}
\end{equation}
其中 $\xi_\ell$ 表示时段、节假日或交通高峰等外生特征；$\mathcal R_{i,\ell-1}^{\mathrm{act}}$ 为上一区间真实到站的随机请求记录，$z_{R,r,i}^{\mathrm{act}}$ 为其实际服务结果。$f_R^{\mathrm{act}}$ 使用真实的到达时刻、SOC 和服务结果，$f_R^{\mathrm{pred}}$ 仅使用时刻 $t_\ell$ 可得的预测量，两者不共享数据字段。由于同一站点的槽位标签没有物理意义，$f_B$ 和 critic 采用可微的置换不变集合编码，actor 的逐槽输出采用相应的置换等变解码；参考 rollout 则按预先固定的对称性消除规则将事件分配给槽位。

\paragraph{Action.}
RL 的动作是各站点、各槽位在当前外层区间的连续请求充电功率：
\begin{equation}
a_\ell=
\{\widehat P_{i,b,\ell}\}_{i\in\mathcal S,b\in\mathcal B_i}.
\label{eq:rl_action}
\end{equation}
actor 先生成无约束逐槽请求，再将其欧氏投影到式\eqref{eq:power_projection}所定义的闭凸集上，并以该唯一投影作为动作 $\widehat P$；该确定映射同时满足单槽功率上限和站级区间充电能量上限。策略概率密度及 PPO 比率均针对这一固定投影参数化所诱导的请求动作分布计算。

\paragraph{Reference rollout.}
为向 MPC 提供未来 $H$ 个区间的请求功率，系统从
$s_{\ell}^{\mathrm{ref}}=s_\ell$ 出发，利用同一 actor 的均值动作和参考状态递推进行策略 rollout：
\begin{equation}
\{\widehat P_{i,b,q}\}_{i,b}
=
\mu_\phi(s_{q}^{\mathrm{ref}}),
\quad q\in\mathcal T_\ell,
\label{eq:rl_power_rollout}
\end{equation}
其中 $\mu_\phi$ 表示策略分布的均值。参考轨迹遵循固定的调用顺序：先按第 3 节规则确定参考路径，再进行 actor rollout，最后由 critic 计算终端价值系数。参考路径的确定只使用求解前已知的位置、SOC、预计到达时间和候选网络；参考执行采用与实际执行相同的预约优先、等待、超时、充电与槽位替换规则，并允许预约软失败。因此，参考状态 $s_q^{\mathrm{ref}}$ 与本轮待求的 $y$、$\alpha$ 和 $z$ 无关，请求序列生成与本轮路径优化之间不存在循环依赖。构造后续参考状态时，区间 $q$ 的预测请求及其名义服务结果进入独立的名义历史编码器 $f_R^{\mathrm{nom}}$；该编码器与 $f_R^{\mathrm{act}}$ 输出维度相同，但不共享输入字段或语义，且只用于参考状态，不写回真实状态。$f_R^{\mathrm{pred}}$ 所需的后续预测窗口由同一需求预测器的条件均值补齐，超出当前 $H$ 步输入的部分采用固定终端预测先验。整条请求功率序列作为本轮 MPC 参数；只有首项请求 $\widehat P_{i,b,\ell}$ 进入真实执行，下一滚动时刻重新生成序列。

\paragraph{Transition.}
给定 $s_\ell$ 和请求功率序列后，系统求解 MPC，得到逐用户最优路径
$\{y_{i,j}^{p,k*}\}$。在途用户中剩余站序发生变化的路径在本轮边界立即发布并结算一次路径调整成本，新站序作为下一轮参考；未到用户的路径只用于当前预测，不对用户发布。为避免内部预测路径进入真实执行，定义边界后的执行路径：在途用户使用本轮已经发布的最优剩余站序，未到用户仍使用日前发布路径 $\bar{\mathbf y}_{A}^{p,k}$；若后者在 $\mathcal I_\ell$ 内实际进入高速公路，仍先按该已发布路径执行，下一边界再作为在途用户更新。记两类用户在 $\mathcal I_\ell$ 内的真实到站记录分别为 $\mathcal H_{A,\ell}^{\mathrm{act}}$ 和 $\mathcal H_{R,\ell}^{\mathrm{act}}$，每条记录携带请求标识、站点、真实到站时刻和真实退回 SOC。

环境随后只在首个外层区间 $\mathcal I_\ell$ 内执行。确定性连续事件算子 $\operatorname{Exec}_\ell$ 的输入仅为真实到站记录、携带等待队列、边界后已发布的执行路径、当前请求功率和观测 SOC；其内部按充电完成、登记到站、队列分配、超时判定的同时刻顺序连续推进：
\begin{equation}
\begin{aligned}
(\mathbf S^{\mathrm{obs}}(t_{\ell+1}),\mathcal W_{\ell+1},
\mathbf s_{A,\ell}^{\mathrm{act}},\mathbf z_{R,\ell}^{\mathrm{act}},
\mathbf f_{A,\ell}^{\mathrm{act}})
=\operatorname{Exec}_\ell\!\Big(&
\mathbf S^{\mathrm{obs}}(t_\ell),\mathcal W_\ell,
\{\widehat P_{i,b,\ell}\}_{i,b},\\
&\{\mathbf y_{A,\ell}^{\mathrm{exec},p,k}\}_{p,k},
\mathcal H_{A,\ell}^{\mathrm{act}},\mathcal H_{R,\ell}^{\mathrm{act}}
\Big).
\end{aligned}
\label{eq:actual_execution_operator}
\end{equation}
其中 $\mathbf y_{A,\ell}^{\mathrm{exec},p,k}=\mathbf y_A^{p,k*}$ 对在途用户成立，而对未到用户固定为 $\bar{\mathbf y}_{A}^{p,k}$。实际执行事件只来自两类真实到站记录和携带队列；预测集合不作为真实到达输入执行算子，预测偏差不触发停止执行，也不要求预测与实际终态强制相等。状态转移可表示为
\begin{equation}
s_{\ell+1}
\sim
\mathbb P\!\left(
\cdot\mid
s_\ell,a_\ell,
\mathrm{MPC}(s_\ell,\widehat{\mathbf P})
\right),
\label{eq:rl_transition}
\end{equation}
其中 $\mathbb P$ 表示随机到达与下一轮预测更新共同诱导的环境状态转移核，$\widehat{\mathbf P}=\{\widehat P_{i,b,q}:i\in\mathcal S,b\in\mathcal B_i,q\in\mathcal T_\ell\}$。这一定义与 MPC 的“优化完整预测域、只执行当前区间”保持一致。

\paragraph{Reward.}
RL reward 使用当前实际执行区间的已实现运营收益，并计入本轮实际发布路径产生的调整成本：
\begin{equation}
r_\ell
=
I_\ell^{\mathrm{act}}
-C_{\mathrm{ch},\ell}
-C_{\mathrm{adj},\ell}
-C_{\mathrm{fail},\ell}.
\label{eq:rl_reward}
\end{equation}
其中 $I_\ell^{\mathrm{act}}$、$C_{\mathrm{ch},\ell}$、$C_{\mathrm{adj},\ell}$、$C_{\mathrm{fail},\ell}$ 分别为区间 $\mathcal I_\ell$ 内真实发生的服务收益、充电成本、路径调整成本和预约超时成本，各项均按真实发生时刻记账，详细展开见附录 D。MPC 目标中的预测随机收益、未来预约超时风险和未来区间成本不进入真实 reward；未到用户的前瞻调整成本也不在当期 reward 中重复实现。该 reward 使 actor 学习逐槽充电功率对当前成本、预约可靠性、路径稳定性和长期电池状态价值的共同影响。


\subsection{RL algorithm}

由于逐槽请求功率是连续动作，且需要在不确定需求环境下保持训练稳定性，本文采用 Proximal Policy Optimization (PPO) 作为 RL 层的训练算法。PPO 属于 actor-critic 方法，actor $\pi_{\phi}(a_t\mid s_t)$ 输出连续请求功率动作的随机策略，critic $V_{\theta}(s_t)$ 估计当前状态的长期价值。选择 PPO 的一个重要原因是其 critic 直接学习 state value function，便于通过式\eqref{eq:marginal_value_soc}提取终端逐槽 SOC 的边际价值并传递给 MPC。

PPO 的 critic 通过最小化价值函数误差进行训练：
\begin{equation}
L_V(\theta)=
\mathbb{E}_t\left[
\left(V_{\theta}(s_t)-\hat{G}_t\right)^2
\right],
\label{eq:value_loss}
\end{equation}
其中 $\hat{G}_t$ 为由采样轨迹计算得到的折扣回报。actor 采用 clipped surrogate objective 更新策略。定义重要性采样比率
\begin{equation}
w_t^{\mathrm{IS}}(\phi)=
\frac{\pi_{\phi}(a_t\mid s_t)}{\pi_{\phi_{old}}(a_t\mid s_t)},
\label{eq:ppo_ratio}
\end{equation}
则策略优化目标为
\begin{equation}
L_{\mathrm{CLIP}}(\phi)=
\mathbb{E}_t\left[
\min\left(
w_t^{\mathrm{IS}}(\phi)\hat{A}_t,
\mathrm{clip}(w_t^{\mathrm{IS}}(\phi),1-\epsilon,1+\epsilon)\hat{A}_t
\right)
\right],
\label{eq:ppo_clip}
\end{equation}
其中 $\hat{A}_t$ 为优势函数估计，$\epsilon$ 为 clip 参数。优势函数由广义优势估计（Generalized Advantage Estimation, GAE）计算：
\begin{equation}
\hat{A}_t=\sum_{n=0}^{\infty}
(\gamma_{\mathrm{RL}}\lambda_{\mathrm{GAE}})^n
e_{t+n}^{\mathrm{TD}},
\quad
e_t^{\mathrm{TD}}=r_t
+\gamma_{\mathrm{RL}}V_{\theta}(s_{t+1})-V_{\theta}(s_t),
\label{eq:gae}
\end{equation}
其中 $\gamma_{\mathrm{RL}}$ 为折扣因子，
$\lambda_{\mathrm{GAE}}$ 为 GAE 参数。外层 RL 决策保持固定步长 $\sigma$，与区间内部的连续物理事件并存。

训练完成后，actor 在每个滚动时刻输出首期逐槽请求功率，并通过式\eqref{eq:rl_power_rollout}生成供 MPC 使用的预测域请求功率序列；实际仅执行首期请求功率。critic 在参考终端状态 $s_{\ell+H}^{\mathrm{ref}}$ 处通过自动微分计算 $V_{\theta}(s)$ 关于每个终端槽位 SOC 分量的梯度，得到 $\lambda^S_{i,b}$；同时按式\eqref{eq:outside_swap_value}评价边界处尚未物理实现交付的价值变化。二者作为固定参数进入线性终端价值项 $\Phi^{\mathrm{RL}}$。由此，RL 层同时向 MPC 提供逐槽请求功率轨迹和长期 SOC 价值，MPC 则反馈剩余路径与连续事件安排，形成闭环协同机制。


\section{Numerical experiments}

% NUMERICAL_EXPERIMENTS_PLACEHOLDER_BEGIN
% 待作者提供实验设置、数据与结果
% NUMERICAL_EXPERIMENTS_PLACEHOLDER_END


\section{Conclusion}

本文针对高速公路换电站网络中预约与随机两类异质用户的运营优化问题，构建了 MPC--RL 分层协同框架。在 MPC 层，本文以滚动时域组织运营：未到预约用户从入口出发、相对日前初始发布路径评价调整，在途预约用户从实时虚拟起点出发、相对最近发布的未执行站序评价调整，站内等待请求跨滚动边界携带且保留原到站时刻与截止时刻；候选网络按 SOC 分档离线生成、在线裁剪，路径调整成本按实际发布结算一次。站内服务以连续时间刻画：请求在到达后最长等待 $\Delta t$，预约优先、类内按到站顺序服务，槽位电池严格在达到 SOC $1$ 后交付，充电在满电时停止、换入退回电池后立即恢复。在 RL 层，actor 输出逐槽请求充电功率作为 MPC 的已知参数，critic 通过终端 SOC 边际价值和边界未决交付的等效价值为 MPC 提供长期价值信息；参考路径先于 actor rollout 确定、随后计算 critic 参数、最后求解 MPC 的调用顺序避免了请求序列与本轮路径决策之间的循环依赖。

本文的模型边界包括：车辆到达下游节点的预计到达时间为外生给定，未建立随机交通模型；换电视为瞬时事件，未区分服务开始与作业结束；终端未决交付价值采用基于同一 critic 的可加近似，其对未来排队与超时风险的刻画是近似的。后续工作可围绕更精细的交通流建模、换电作业时长以及终端价值的精确刻画展开，并通过系统的数值实验评估框架在不同需求强度、电价模式和库存配置下的表现。


\appendix

\section{Candidate Events and Effective Arrival Chains}
\label{app:events}

本附录给出 MPC 中候选换电事件的构造、激活量与有效到达链。站内服务使用通用请求索引 $r$；在需要区分来源时，预约候选事件记为 $\varepsilon=(p,k,j,i)$，由入弧 $(j,i)$ 唯一确定退回 SOC。

站点 $i$、区间 $q\in\mathcal T_\ell$ 的候选预约事件集合为
\begin{equation}
\mathcal E_{i,q}^{A}
=
\left\{(p,k,j,i):p\in\mathcal P,\ k\in\mathcal K_\ell^p,
\ (j,i)\in\mathcal A^{p,k},
\ q(t_{A,i}^{p,k})=q\right\},
\label{eq:swap_event_sets}
\end{equation}
其中 $t_{A,i}^{p,k}$ 是求解前已知的连续预计到达时间：未到用户由预约入口时刻与行程时间计算，在途用户由实时位置外生给定。该集合由已知信息枚举，不以未知路径变量 $y=1$ 为成员条件。上一滚动时刻遗留下来、仍未超时且未服务的等待请求构成携带事件集合 $\mathcal E_{i,\ell}^{A,\mathrm{wait}}$（预约）与 $\mathcal E_{i,\ell}^{R,\mathrm{wait}}$（随机），它们是既有请求记录，保留原请求标识、原到站时刻与原截止时刻，不因原到达区间不属于当前 $\mathcal T_\ell$ 而被删除。随机候选集合为 $\mathcal E_{i,q}^{R}=\mathcal R_{i,q}$。区间 $q$ 的新到请求集仅为
$\mathcal E_{i,q}=\mathcal E_{i,q}^{A}\cup\mathcal E_{i,q}^{R}$；
携带请求在站级集合中只出现一次：
\begin{equation}
\begin{aligned}
\mathcal E_i^A
&=\left(\bigcup_{q\in\mathcal T_\ell}\mathcal E_{i,q}^A\right)
\cup\mathcal E_{i,\ell}^{A,\mathrm{wait}},\\
\mathcal E_i^R
&=\left(\bigcup_{q\in\mathcal T_\ell}\mathcal E_{i,q}^R\right)
\cup\mathcal E_{i,\ell}^{R,\mathrm{wait}},
\qquad
\mathcal E_i=\mathcal E_i^A\cup\mathcal E_i^R.
\end{aligned}
\label{eq:station_event_union}
\end{equation}

每个候选预约事件的路径激活量与退回 SOC 为
\begin{equation}
(\delta_\varepsilon,\rho_\varepsilon)=
\begin{cases}
(y_{j,i}^{p,k},\rho_{A,j,i}^{p,k}),
&\varepsilon=(p,k,j,i)\in\mathcal E_{i,q}^{A},\\
(1,\rho_\varepsilon^{\mathrm{rec}}),
&\varepsilon\in\mathcal E_{i,\ell}^{A,\mathrm{wait}},
\end{cases}
\label{eq:event_activation_and_soc}
\end{equation}
其中 $\rho_{A,j,i}^{p,k}$ 按式\eqref{eq:return_soc}计算，携带事件的 $\rho_\varepsilon^{\mathrm{rec}}$ 为其请求记录中的已知退回 SOC。随机请求的存在量固定为 $1$，是否服务由 $z_r$ 控制。

为避免用户在上游超时后仍在下游重复服务或重复计罚，定义预约事件的有效到达量 $a_\varepsilon^A\in\{0,1\}$。对以本轮起点 $o_\ell^{p,k}$ 为首段入弧的首站事件，令
$a_\varepsilon^A=y_{j,i}^{p,k}$（$j=o_\ell^{p,k}$）；
对携带等待事件，令 $a_\varepsilon^A=1$；
对其余沿弧 $(j,i)$ 从上游站 $j$ 到达的事件，令
\begin{subequations}
\label{eq:reservation_alive_chain}
\begin{align}
a_\varepsilon^A&\le y_{j,i}^{p,k},\\
a_\varepsilon^A&\le
\sum_{\varepsilon'\in\mathcal E^A(p,k,j)}s_{\varepsilon'}^A,\\
a_\varepsilon^A&\ge y_{j,i}^{p,k}
+\sum_{\varepsilon'\in\mathcal E^A(p,k,j)}s_{\varepsilon'}^A-1,
\end{align}
\end{subequations}
其中 $\mathcal E^A(p,k,j)$ 为同一用户到达上游站 $j$ 的候选预约事件集合，所选路径上至多有一个对应服务量为 $1$。因此，一次上游超时使该用户的下游预约事件失活，且只结算真实发生的一次失败。已执行路径前缀上的服务以观测事实作为参数继承，不在本轮重复求解。未被激活或已失活的事件满足 $s_\varepsilon^A=f_\varepsilon^A=\omega_\varepsilon^A=0$，不会被误计为超时。


\section{Waiting, Priority, and Timeout Constraints}
\label{app:waiting}

令 $T=t_{\ell+H}$，并以 $\mathcal E_i=\mathcal E_i^A\cup\mathcal E_i^R$ 统称站点 $i$ 在本轮可能进入队列的预约、随机和携带请求。对任一请求 $r$，定义存在量 $e_r$、服务量 $\varsigma_r$、域内超时或流失量 $\chi_r$ 和边界未决量 $\omega_r$：
\[
(e_r,\varsigma_r,\chi_r,\omega_r)
=
\begin{cases}
(a_\varepsilon^A,s_\varepsilon^A,f_\varepsilon^A,\omega_\varepsilon^A),
&r=\varepsilon\text{ 为预约事件},\\
(1,z_r,1-z_r-\omega_r^R,\omega_r^R),
&r\text{ 为随机请求}.
\end{cases}
\]
令 $h_r^T=\mathbb I\{t_r^{\mathrm{ddl}}\ge T\}$ 为求解前已知的跨界标记。预约事件和随机请求分别满足
\begin{subequations}
\label{eq:request_outcome_conservation}
\begin{align}
s_\varepsilon^A+f_\varepsilon^A+\omega_\varepsilon^A
&=a_\varepsilon^A,&
\omega_\varepsilon^A&\le h_\varepsilon^T,&
f_\varepsilon^A&\le1-h_\varepsilon^T,\label{eq:reservation_outcome_boundary}\\
z_r+\omega_r^R&\le1,&
\omega_r^R&\le h_r^T,&
z_r+\omega_r^R&\ge h_r^T.\label{eq:random_outcome_boundary}
\end{align}
\end{subequations}
因此，$t_r^{\mathrm{ddl}}<T$ 时有 $\omega_r=0$，域内未服务请求只能在截止时刻超时或流失；$t_r^{\mathrm{ddl}}\ge T$ 时有 $\chi_r=0$，随机请求还由最后一个不等式强制为“域内服务”或“跨界未决”二者之一，不能在完整等待窗结束前消失。等号情形 $t_r^{\mathrm{ddl}}=T$ 归入未决，并由下一轮在时刻 $T$ 处理。




附录 C 使用固定位置集
$\mathcal M_i=\{0,\ldots,\overline M_i-1\}$，
并定义 $\vartheta_{i,m}=1$ 当且仅当位置 $m$ 位于半开预测域内；终端及填充位置取 $\vartheta_{i,m}=0$。对本轮首次到站的请求，二元标记 $\iota_{i,r,m}^{\mathrm{arr}}$ 将其到站时刻指派给唯一活动位置；对 $t_r^{\mathrm{ddl}}<T$ 的请求，$\iota_{i,r,m}^{\mathrm{ddl}}$ 类似地指派截止时刻。使用平移时间
$\bar\tau_{i,m}=\tau_{i,m}-t_\ell$、
$\bar t=t-t_\ell$ 和合法界
$M_T=T-t_\ell=H\sigma$，有
\begin{subequations}
\label{eq:event_time_assignment}
\begin{align}
-M_T(1-\iota_{i,r,m}^{\mathrm{arr}})
&\le \bar\tau_{i,m}-\bar t_r^{\mathrm{arr}}
\le M_T(1-\iota_{i,r,m}^{\mathrm{arr}}),\\
-M_T(1-\iota_{i,r,m}^{\mathrm{ddl}})
&\le \bar\tau_{i,m}-\bar t_r^{\mathrm{ddl}}
\le M_T(1-\iota_{i,r,m}^{\mathrm{ddl}}),\\
\sum_{m\in\mathcal M_i}\iota_{i,r,m}^{\mathrm{arr}}&=1,
\quad t_r^{\mathrm{arr}}\in[t_\ell,T),&
\sum_{m\in\mathcal M_i}\iota_{i,r,m}^{\mathrm{ddl}}&=1,
\quad t_r^{\mathrm{ddl}}\in[t_\ell,T),\\
\iota_{i,r,m}^{\mathrm{arr}}&\le \vartheta_{i,m},&
\iota_{i,r,m}^{\mathrm{ddl}}&\le \vartheta_{i,m}.
\end{align}
\end{subequations}
候选到站时刻即使对应预约事件未激活也保留，故固定位置上界不依赖未知路径变量；同一物理位置可同时携带多个到站或截止标签。携带请求不重复登记，而在位置 $0$ 以前初始化。

令 $w_{i,r,m}^-$、$w_{i,r,m}^{\mathrm{arr}}$、
$w_{i,r,m}^{\mathrm{sw}}$ 和 $w_{i,r,m}^+$ 分别表示请求在位置 $m$ 的登记前、登记后、分配后和超时判定后的等待状态；$x_{i,r,m}=1$ 表示请求在该位置获得服务。新到请求的有效登记量
$a_{i,r,m}^{\mathrm{in}}=e_r\iota_{i,r,m}^{\mathrm{arr}}$
采用三个标准二元乘积包络。对所有 $m\in\mathcal M_i$，队列递推为
\begin{subequations}
\label{eq:queue_event_flow}
\begin{align}
w_{i,r,m}^{\mathrm{arr}}
&=w_{i,r,m}^-+a_{i,r,m}^{\mathrm{in}},&
x_{i,r,m}&\le w_{i,r,m}^{\mathrm{arr}},&
x_{i,r,m}&\le \vartheta_{i,m},\\
w_{i,r,m}^{\mathrm{sw}}
&=w_{i,r,m}^{\mathrm{arr}}-x_{i,r,m},&
\chi_{i,r,m}^{\mathrm{to}}
&=\iota_{i,r,m}^{\mathrm{ddl}}w_{i,r,m}^{\mathrm{sw}},\\
w_{i,r,m}^{+}
&=w_{i,r,m}^{\mathrm{sw}}-\chi_{i,r,m}^{\mathrm{to}},&
w_{i,r,m+1}^{-}&=w_{i,r,m}^{+}.
\end{align}
\end{subequations}
其中两个二元乘积均使用
$y\le y_1$、$y\le y_2$、
$y\ge y_1+y_2-1$ 的线性包络。携带请求取
$w_{i,r,0}^-=e_r$，本轮新到请求取
$w_{i,r,0}^-=0$；两类集合互斥。终端和填充位置没有到站或截止标签，且 $\vartheta_{i,m}=0$ 强制 $x_{i,r,m}=0$，上述递推只原样携带队列。结果变量满足
\begin{equation}
\varsigma_r=\sum_{m\in\mathcal M_i}x_{i,r,m},
\qquad
\chi_r=\sum_{m\in\mathcal M_i}\chi_{i,r,m}^{\mathrm{to}},
\qquad
\omega_r=w_{i,r,\overline M_i}^{-}.
\label{eq:queue_outcome_link}
\end{equation}

站内请求使用固定全序
\begin{equation}
r'\prec_i r
\Longleftrightarrow
\bigl(\mathbb I\{r'\in\mathcal E_i^R\},
t_{r'}^{\mathrm{arr}},\operatorname{id}(r')\bigr)
<_{\mathrm{lex}}
\bigl(\mathbb I\{r\in\mathcal E_i^R\},
t_r^{\mathrm{arr}},\operatorname{id}(r)\bigr),
\label{eq:reservation_priority}
\end{equation}
即预约在前，两类内部按原到站时刻和固定编号排序。每个位置施加
\begin{equation}
x_{i,r,m}
\le x_{i,r',m}+1-w_{i,r',m}^{\mathrm{arr}},
\qquad r'\prec_i r,\quad m\in\mathcal M_i.
\label{eq:reservation_service_order}
\end{equation}
若高优先请求尚未到站或已离队，则
$w_{i,r',m}^{\mathrm{arr}}=0$，不会为其扣留电池；若其正在等待，低优先请求只有在高优先请求同刻也获服务时才可获服务。再令附录 C 中的 $h_{i,b,m}^+$ 表示分配后槽位是否仍满电，施加功守恒约束
\begin{equation}
w_{i,r,m}^{\mathrm{sw}}+h_{i,b,m}^{+}
\le1+(1-\vartheta_{i,m}),
\qquad r\in\mathcal E_i,\ b\in\mathcal B_i,\ m\in\mathcal M_i.
\label{eq:work_conservation}
\end{equation}
活动位置上该式排除“仍有合格等待请求却留下未用满电槽位”的解；在终端或填充位置，右端的松弛项关闭分配要求。与附录 C 的匹配容量结合后，每个活动位置恰好服务优先队列中能够由当前满电槽位覆盖的最长前缀。

令 $\bar t_r^{\mathrm{sw}}=t_r^{\mathrm{sw}}-t_\ell$。服务时刻由
\begin{subequations}
\label{eq:service_window}
\begin{align}
-M_T(1-x_{i,r,m})
&\le \bar t_r^{\mathrm{sw}}-\bar\tau_{i,m}
\le M_T(1-x_{i,r,m}),\\
\bar t_r^{\mathrm{sw}}
&\ge \max\{0,\bar t_r^{\mathrm{arr}}\}
-M_T(1-\varsigma_r),\\
\bar t_r^{\mathrm{sw}}
&\le \min\{M_T,\bar t_r^{\mathrm{ddl}}\}
+M_T(1-\varsigma_r),\qquad
0\le\bar t_r^{\mathrm{sw}}\le M_T
\end{align}
\end{subequations}
确定，其中第一式对 $m\in\mathcal M_i$ 施加。由于 $x_{i,r,m}\le \vartheta_{i,m}$，成功服务严格发生在 $T$ 以前；若服务位置同时带有截止标签，式\eqref{eq:queue_event_flow}先扣除服务再计算超时。

附录 C 的区间归属量 $\zeta_{i,m,q}=1$ 精确表示
$\tau_{i,m}\in[t_q,t_{q+1})$。为线性化
$x_{i,r,m}\zeta_{i,m,q}$，引入
$\Gamma_{i,r,m,q}\in\{0,1\}$ 并施加
\[
\Gamma_{i,r,m,q}\le x_{i,r,m},\qquad
\Gamma_{i,r,m,q}\le\zeta_{i,m,q},\qquad
\Gamma_{i,r,m,q}\ge x_{i,r,m}+\zeta_{i,m,q}-1.
\]
服务价格区间由
\begin{equation}
\gamma_{i,r,q}
=\sum_{m\in\mathcal M_i}\Gamma_{i,r,m,q},
\qquad
\sum_{q\in\mathcal T_\ell}\gamma_{i,r,q}=\varsigma_r
\label{eq:service_price_bin}
\end{equation}
唯一确定。边界累计定义保证恰在 $t_q$ 的服务归入新区间 $q$，而 $T$ 及填充位置满足 $\sum_q\zeta_{i,m,q}=0$，不属于任何本轮价格区间。

本附录中的存在、结果、时刻指派、队列状态、事件服务量和价格乘积量均为二元量；$t_r^{\mathrm{sw}}$ 与 $\tau_{i,m}$ 为连续时刻。所有 Big-$M$ 关系均在平移后的预测域 $[0,H\sigma]$ 上使用 $M_T=H\sigma$，不依赖运营日的绝对时间原点。

\section{Continuous Charging and Slot Matching}
\label{app:charging}

先构造与求解变量无关的固定位置上界。令
\[
\mathcal F_i=\{t_\ell,t_{\ell+1},\ldots,t_{\ell+H}\}
\cup\{t_r^{\mathrm{arr}}:
r\in\mathcal E_i,\ t_r^{\mathrm{arr}}\in[t_\ell,T)\}
\cup\{t_r^{\mathrm{ddl}}:
r\in\mathcal E_i,\ t_r^{\mathrm{ddl}}\in[t_\ell,T)\}
\]
为去重后的外层边界、候选到站和域内截止时刻集合。每个槽位在初始状态后至多首次充满一次，每次服务后至多再次充满一次，因此
\[
\overline M_i=|\mathcal F_i|+|\mathcal B_i|+|\mathcal E_i|,
\qquad
\overline{\mathcal M}_i=\{0,\ldots,\overline M_i\},
\qquad
\mathcal M_i=\{0,\ldots,\overline M_i-1\}
\]
给出有效的有限上界。实际物理时刻按升序装入
$\tau_{i,0},\ldots,\tau_{i,\overline M_i}$；同刻的边界、到站、截止和一个或多个充满事件合并为一个位置，实际时刻用完后的槽位全部以 $T$ 填充。于是
\begin{subequations}
\label{eq:event_position_grid}
\begin{align}
t_\ell=\tau_{i,0}
&\le\tau_{i,1}\le\cdots
\le\tau_{i,\overline M_i}=T,\\
d_{i,m}&=\tau_{i,m+1}-\tau_{i,m},&
0\le d_{i,m}&\le\sigma\vartheta_{i,m},
\qquad m\in\mathcal M_i,
\end{align}
\end{subequations}
其中 $\vartheta_{i,m}=1$ 当且仅当 $\tau_{i,m}<T$；填充位置取 $\vartheta_{i,m}=0$ 并由上式强制 $d_{i,m}=0$。

为精确处理半开区间，令
$\iota_{i,q,m}^{\mathrm{bd}}=1$ 表示边界 $t_q$ 位于位置 $m$。对
$q\in\{\ell,\ldots,\ell+H\}$ 和
$m\in\overline{\mathcal M}_i$，使用
\begin{subequations}
\label{eq:boundary_position_assignment}
\begin{align}
-M_T(1-\iota_{i,q,m}^{\mathrm{bd}})
&\le\bar\tau_{i,m}-\bar t_q
\le M_T(1-\iota_{i,q,m}^{\mathrm{bd}}),\\
\sum_{m\in\overline{\mathcal M}_i}
\iota_{i,q,m}^{\mathrm{bd}}&=1,
&\iota_{i,\ell,0}^{\mathrm{bd}}&=1.
\end{align}
\end{subequations}
边界标签按上述合并规则放在该时刻的第一个位置。定义
\begin{equation}
B_{i,q,m}=\sum_{n=0}^{m}\iota_{i,q,n}^{\mathrm{bd}},
\qquad
\zeta_{i,m,q}=B_{i,q,m}-B_{i,q+1,m},
\qquad
\vartheta_{i,m}=\sum_{q\in\mathcal T_\ell}\zeta_{i,m,q}.
\label{eq:boundary_interval_affiliation}
\end{equation}
由此，$\zeta_{i,m,q}=1$ 当且仅当
$\tau_{i,m}\in[t_q,t_{q+1})$；恰在 $t_q$ 的位置归入新区间 $q$，而 $T$ 及其填充位置不属于本轮任何区间。物理时刻严格递增直至第一次出现 $T$，相同时刻只保留一个位置；该索引规范化是事件递推的定义，不借助数值 $\epsilon$ 区分同刻事件。

槽位 $b$ 在位置 $m$ 分配前后的 SOC 分别记为
$S_{i,b,m}^{-}$ 和 $S_{i,b,m}^{+}$，相应服务就绪标记为
$h_{i,b,m}^{-}$ 和 $h_{i,b,m}^{+}$。本轮初始条件为
\begin{equation}
S_{i,b,0}^{-}=S_{i,b}^{\mathrm{obs}}(t_\ell),
\qquad
h_{i,b,0}^{-}=\mathbb I\{S_{i,b}^{\mathrm{obs}}(t_\ell)=1\}.
\label{eq:initial_battery_soc}
\end{equation}
对活动位置 $m$，下一个固定候选时刻是
\[
\Theta_{i,m}^{\mathrm{fix}}
=\min\{t\in\mathcal F_i:t>\tau_{i,m}\}.
\]
若 $\zeta_{i,m,q}=1$、$h_{i,b,m}^{+}=0$ 且
$\widehat P_{i,b,q}>0$，槽位的理论充满候选为
\begin{equation}
\Theta_{i,b,m,q}^{\mathrm{full}}
=\tau_{i,m}
+\frac{E_B(1-S_{i,b,m}^{+})}
{\eta_i\widehat P_{i,b,q}}.
\label{eq:full_time_candidate}
\end{equation}
仅当该时刻严格位于 $(\tau_{i,m},T)$ 时才将其加入候选；等于 $T$ 的充满由下一轮处理。令 $\mathcal C_{i,m}$ 包含上述下一个固定候选以及所有有效充满候选，并将候选 $c$ 的时刻记为 $\Theta_{i,c,m}$。精确最小时刻标签定义为
\[
\psi_{i,c,m}
=\mathbb I\!\left\{
\Theta_{i,c,m}=\min_{c'\in\mathcal C_{i,m}}
\Theta_{i,c',m}\right\}.
\]
即并列最小的标签全部取 $1$，而非任意选择其中一个；$\psi$ 是确定性事件生成器输出的派生标签，不是可自由取值的决策变量。令
$\bar\Theta_{i,c,m}=\Theta_{i,c,m}-t_\ell$。
在给定该标签后，与连续时刻连接的 indicator/Big-$M$ 接口为
\begin{subequations}
\label{eq:next_physical_event}
\begin{align}
\bar\tau_{i,m+1}&\le\bar\Theta_{i,c,m},
&&c\in\mathcal C_{i,m},\\
\bar\tau_{i,m+1}&\ge
\bar\Theta_{i,c,m}-M_T(1-\psi_{i,c,m}),
&&c\in\mathcal C_{i,m},\\
\sum_{c\in\mathcal C_{i,m}}\psi_{i,c,m}&\ge1,
&&\vartheta_{i,m}=1.
\end{align}
\end{subequations}
候选有效性是 $\zeta_{i,m,q}$、$1-h_{i,b,m}^{+}$ 和正功率条件的逻辑与，可用标准二元乘积包络实现。对每个槽位，令
\[
\varpi_{i,b,m}^{\mathrm{full}}
=\sum_{\substack{c\in\mathcal C_{i,m}:\\
c\text{ 为槽位 }b\text{ 的充满候选}}}
\psi_{i,c,m},
\qquad
g_{i,b,m}=h_{i,b,m}^{+}+\varpi_{i,b,m}^{\mathrm{full}},
\qquad
h_{i,b,m+1}^{-}=g_{i,b,m},
\]
并施加 $\varpi_{i,b,m}^{\mathrm{full}}\le1-h_{i,b,m}^{+}$。这样，充满与固定事件同刻时，充满标签不会因另一个并列候选被选中而丢失；$g$ 也不再是可任意关闭的“$S=1$”分支。每个物理位置按“前段充满完成 $\rightarrow$ 外层边界功率切换 $\rightarrow$ 到站登记 $\rightarrow$ 队列分配 $\rightarrow$ 未服务请求超时”的顺序处理。终点 $T$ 只保留事件前快照，不执行这些阶段。

以 $u_{i,b,m}$ 表示槽位在事件段内因段首已经满电而停止充电的时长。由于式\eqref{eq:next_physical_event}不会让未满槽位越过首次充满时刻，$u$ 恰为
$d_{i,m}h_{i,b,m}^{+}$，其线性化为
\begin{subequations}
\label{eq:full_stop_indicator}
\begin{align}
0\le u_{i,b,m}&\le d_{i,m},&
u_{i,b,m}&\le\sigma h_{i,b,m}^{+},\\
u_{i,b,m}&\ge d_{i,m}
-\sigma(1-h_{i,b,m}^{+}).
\end{align}
\end{subequations}
再令 $\upsilon_{i,b,m,q}^{\mathrm{ch}}$ 为事件段 $m$ 在区间 $q$ 内的实际充电时长。乘积
$\zeta_{i,m,q}(d_{i,m}-u_{i,b,m})$ 由
\begin{subequations}
\label{eq:interval_charging_duration}
\begin{align}
0\le\upsilon_{i,b,m,q}^{\mathrm{ch}}
&\le\sigma\zeta_{i,m,q},\\
\upsilon_{i,b,m,q}^{\mathrm{ch}}
&\le d_{i,m}-u_{i,b,m},\\
\upsilon_{i,b,m,q}^{\mathrm{ch}}
&\ge d_{i,m}-u_{i,b,m}
-\sigma(1-\zeta_{i,m,q}),\\
\sum_{q\in\mathcal T_\ell}
\upsilon_{i,b,m,q}^{\mathrm{ch}}
&=d_{i,m}-u_{i,b,m}
\end{align}
\end{subequations}
精确给出。SOC 递推于是写为
\begin{equation}
S_{i,b,m+1}^{-}
=S_{i,b,m}^{+}
+\frac{\eta_i}{E_B}
\sum_{q\in\mathcal T_\ell}
\widehat P_{i,b,q}
\upsilon_{i,b,m,q}^{\mathrm{ch}},
\qquad 0\le S_{i,b,m+1}^{-}\le1.
\label{eq:appendix_charging_recursion}
\end{equation}
请求功率为 $0$ 时，未满槽位保持原 SOC 且不生成伪充满事件；段首已满的槽位取 $u=d$。槽位在 $T$ 恰好充满时，式\eqref{eq:appendix_charging_recursion}仍给出 $S(T^-)=1$，但终点不生成可服务事件。

对每个 $m\in\mathcal M_i$，引入
$\alpha_{i,b,r,m}\in\{0,1\}$；
$\alpha_{i,b,r,m}=1$ 表示请求 $r$ 在该位置由槽位 $b$ 服务。令
$a_{i,b,m}=\sum_{r\in\mathcal E_i}\alpha_{i,b,r,m}$。
附录 B 的服务量与槽位匹配满足
\begin{subequations}
\label{eq:battery_event_assignment}
\begin{align}
\sum_{b\in\mathcal B_i}\alpha_{i,b,r,m}
&=x_{i,r,m},
&&r\in\mathcal E_i,\ m\in\mathcal M_i,\\
\sum_{r\in\mathcal E_i}\alpha_{i,b,r,m}
&=a_{i,b,m},
&&b\in\mathcal B_i,\ m\in\mathcal M_i,\\
a_{i,b,m}&\le h_{i,b,m}^{-},&
a_{i,b,m}&\le \vartheta_{i,m}.
\end{align}
\end{subequations}
因此一个请求恰好匹配一个槽位，一个槽位同刻至多服务一个请求，且只有 SOC 精确为 $1$ 的就绪槽位可交付。本文的有效换电请求均满足 $0\le\rho_r<1$。
\begin{subequations}
\label{eq:continuous_swap_transition}
\begin{align}
h_{i,b,m}^{+}
&=h_{i,b,m}^{-}-a_{i,b,m},\\
S_{i,b,m}^{+}
&=S_{i,b,m}^{-}-a_{i,b,m}
+\sum_{r\in\mathcal E_i}\rho_r\alpha_{i,b,r,m}.
\end{align}
\end{subequations}
式\eqref{eq:continuous_swap_transition}给出分配后的状态。若槽位未被使用，前后状态相同；若被使用，则
$S_{i,b,m}^{-}=1$，服务后的 SOC 恰为退回值 $\rho_r$，并立即按位置所属区间的同一请求功率重新充电。

槽位编号按固定升序排列。对 $b'<b$，以及同站优先序 $r'\prec_i r$，施加
\begin{subequations}
\label{eq:smallest_ready_slot}
\begin{align}
a_{i,b,m}
&\le a_{i,b',m}+1-h_{i,b',m}^{-},\\
\alpha_{i,b,r',m}+\alpha_{i,b',r,m}
&\le1.
\end{align}
\end{subequations}
第一式保证使用较大编号槽位时，所有更小且已就绪的槽位均已使用；第二式保证同刻并行服务时，高优先请求不会获得比低优先请求更大的槽位编号。结合式\eqref{eq:work_conservation}，分配结果唯一遵循“优先队列最长可服务前缀—最小就绪槽位”的规则。同一槽位可在同一外层区间内多次服务，但式\eqref{eq:battery_event_assignment}要求它在每次服务前重新出现 $h_{i,b,m}^{-}=1$，排除了未真实充满便再次交付。

槽位在区间 $q$ 的实际购电量及站级限额为
\begin{subequations}
\label{eq:event_interval_energy}
\begin{align}
E_{i,b,q}^{\mathrm{act}}
&=\sum_{m\in\mathcal M_i}
\widehat P_{i,b,q}
\upsilon_{i,b,m,q}^{\mathrm{ch}},\\
\sum_{b\in\mathcal B_i}E_{i,b,q}^{\mathrm{act}}
&\le\overline E_{i,q},
\qquad i\in\mathcal S,\ q\in\mathcal T_\ell.
\end{align}
\end{subequations}
于是式\eqref{eq:chargingcost}等价于
$C_{\mathrm{ch}}=\sum_{i,q,b}e_{i,q}E_{i,b,q}^{\mathrm{act}}$；
边界位置把充电段严格切开，不会沿用旧区间功率或电价。终端 SOC 使用
$S_{i,b}(T^-)=S_{i,b,\overline M_i}^{-}$；终端及填充位置的 $\vartheta=0$，因而不存在服务或超时。变量域为
\[
\begin{gathered}
0\le S_{i,b,m}^{-},S_{i,b,m}^{+}\le1,
\qquad d_{i,m},u_{i,b,m},
\upsilon_{i,b,m,q}^{\mathrm{ch}}\ge0,\\
h_{i,b,m}^{-},h_{i,b,m}^{+},g_{i,b,m},
\varpi_{i,b,m}^{\mathrm{full}},a_{i,b,m},\vartheta_{i,m},
\zeta_{i,m,q}\in\{0,1\},\\
\alpha_{i,b,r,m},\iota_{i,q,m}^{\mathrm{bd}}
\in\{0,1\}.
\end{gathered}
\]
其中候选有效性、精确等时合并和 $\psi$ 的全并列规范均由确定性事件生成器给出；式\eqref{eq:next_physical_event}只连接其输出与连续事件时刻，不能另行关闭某个并列充满标签。完整实现因而是“确定性规范化事件生成器 + 固定维度线性条件模型”的组合：全并列的反向蕴含由事件生成器精确处理，而不以带数值 $\epsilon$ 的标准 Big-$M$ 约束近似；在此规范下，持续时间、SOC、能量和二元乘积均由上述线性关系展开。



\section{Publication, Boundary State, and Accounting}
\label{app:accounting}

本附录展开路径发布的参考对齐、跨边界状态继承与真实记账规则。

\textbf{参考对齐与发布。}式\eqref{eq:reference_visit_indicator}中的参考访问量按用户状态区分：未到用户的参考为日前初始发布路径的站点访问量
$\bar x_{A,i}^{p,k}=\sum_{\{j\mid(j,i)\in\bar{\mathcal A}^{p,k}\}}\bar y_{A,j,i}^{p,k}$；
在途用户的参考为最近一次实际发布路径 $y_A^{\mathrm{pub},p,k}$ 的未执行站序对应的站点访问量 $x_{A,i\mid\ell}^{\mathrm{pub},p,k}$。比较范围为当前候选实体站点与参考未执行站点的并集，缺失访问量补零；已执行路径前缀不参与比较。在途用户的剩余站序在本轮发生变化时，于本轮边界立即发布，结算一次路径调整成本，并将新站序写入 $y_A^{\mathrm{pub},p,k}$ 作为下一轮参考；虚拟起点随滚动前移而剩余站序不变时，$d_{p,k}=0$，不产生路径调整成本。未到用户的内部优化路径不发布、不覆盖 $\bar y_A$，也不进入当期实际记账。

\textbf{边界状态。}滚动边界处继承的状态包括：逐槽 SOC $S_{i,b}^{\mathrm{obs}}(t_{\ell+1})$（含正在充电的槽位）、仍未超时的等待队列 $\mathcal W_{\ell+1}$（保留原请求标识、原到站时刻与原截止时刻）以及已发布剩余路径。对预约用户 $(p,k)$ 定义边界存活变量 $A_{p,k}^{\mathrm{bd}}\in\{0,1\}$；若其任一域内激活事件超时，则后续交付全部失活。令 $\mathcal E_{p,k}^{A,\mathrm{in}}$ 为该用户在预测边界前可能到达的固定候选事件集合，则
\begin{subequations}
\label{eq:boundary_alive}
\begin{align}
A_{p,k}^{\mathrm{bd}}
&\ge 1-\sum_{\varepsilon\in\mathcal E_{p,k}^{A,\mathrm{in}}}
f_\varepsilon^A,\\
A_{p,k}^{\mathrm{bd}}
&\le 1-f_\varepsilon^A,
\qquad \varepsilon\in\mathcal E_{p,k}^{A,\mathrm{in}},
\end{align}
\end{subequations}
即 $A_{p,k}^{\mathrm{bd}}=1$ 当且仅当该用户没有域内激活事件超时；若域内事件集为空，则固定 $A_{p,k}^{\mathrm{bd}}=1$。正文中的 $\overline{\mathcal D}_{\ell}^{A}$ 在求解前枚举所有尚可能影响终端库存的候选预约交付。对记录 $h\in\overline{\mathcal D}_{\ell}^{A}$，属性 $i(h)$、$p(h)$、$k(h)$ 和 $\rho_h$ 分别表示交付站点、用户 O--D 对、用户编号和退回 SOC。令 $\delta_h$ 为其已知或路径诱导的激活量：普通候选交付取相应入弧变量，携带请求取 $1$；令 $s_h^{\mathrm{in}}$ 为其域内服务量，若该交付在边界前不可能服务则固定为 $0$。未决交付是 $\delta_h$、用户边界存活和未在域内完成三者的逻辑与：
\begin{subequations}
\label{eq:pending_delivery_alive}
\begin{align}
w_h^{\mathrm{pend}}&\le\delta_h,&
w_h^{\mathrm{pend}}&\le A_{p(h),k(h)}^{\mathrm{bd}},\\
w_h^{\mathrm{pend}}&\le1-s_h^{\mathrm{in}},&
w_h^{\mathrm{pend}}
&\ge\delta_h+A_{p(h),k(h)}^{\mathrm{bd}}
-s_h^{\mathrm{in}}-1.
\end{align}
\end{subequations}
因此，预计在边界后到达的交付以及边界处仍在等待的域内到达请求均可取 $w_h^{\mathrm{pend}}=1$；若服务已经在域内完成，则 $s_h^{\mathrm{in}}=1$ 强制其未决指示为 $0$。同一交付不会同时进入已实现终端 SOC 和未决价值，式\eqref{eq:rl_terminal_value}由此覆盖全部尚未物理实现的预约交付，而不再仅按预计到达时刻是否位于域外分类。

\textbf{实际记账。}各成本与收益按真实发生时刻结算，且跨轮只结算一次。
\begin{align}
I_\ell^{A,\mathrm{act}}
&=E_B\sum_{i\in\mathcal S}
\sum_{\substack{r:\ \text{预约},\ i(r)=i,\\
t_r^{\mathrm{sw,act}}\in\mathcal I_\ell}}
\pi_{i,\ell}^{\mathrm{sw}}
(1-\rho_r^{\mathrm{act}}),
\label{eq:realized_reservation_income}\\
I_\ell^{R,\mathrm{act}}
&=E_B\sum_{i\in\mathcal S}
\sum_{\substack{r\in\mathcal R_{i}^{\mathrm{act}}:\\
t_r^{\mathrm{sw,act}}\in\mathcal I_\ell}}
\pi_{i,\ell}^{\mathrm{sw}}
(1-\mathrm{soc}_R^{i,r}),
\label{eq:realized_random_income}\\
C_{\mathrm{ch},\ell}
&=\sum_{i\in\mathcal S}\sum_{b\in\mathcal B_i}
e_{i,\ell}
\int_{\mathcal I_\ell}P_{i,b}^{\mathrm{act}}(t)\,\mathrm dt,
\label{eq:realized_charging_cost}\\
C_{\mathrm{adj},\ell}
&=\kappa\sum_{p\in\mathcal P}
\sum_{k\in\mathcal K_\ell^{\mathrm{enr},p}}
\mathbb I\!\left\{
\text{用户 }k\text{ 的剩余站序在本轮发布时发生变化}
\right\},
\label{eq:realized_adjustment_cost}\\
C_{\mathrm{fail},\ell}
&=c^{\mathrm{fail}}\sum_{i\in\mathcal S}
\sum_{\substack{r:\ \text{预约},\ i(r)=i,\\
t_r^{\mathrm{ddl}}\in\mathcal I_\ell}}
f_r^{\mathrm{act}},
\label{eq:realized_reservation_failure_cost}
\end{align}
上述各式给出区间 $\mathcal I_\ell$ 的已实现量。其中 $t_r^{\mathrm{sw,act}}$、$\rho_r^{\mathrm{act}}$ 和 $f_r^{\mathrm{act}}$ 为执行算子记录的真实服务时刻、真实退回 SOC 与真实超时结果。服务收益按实际服务时刻所在价格区间结算，充电成本按区间内实际功率积分结算，路径调整成本按实际发布时刻结算，预约失败成本按截止时刻所在执行区间结算；上期到站、本期服务的请求只在实际服务区间结算一次收入。式\eqref{eq:rl_reward}中的各项即由上述已实现量构成。成功服务请求的已实现等待时间为 $t_r^{\mathrm{sw,act}}-t_r^{\mathrm{arr}}$；在截止时刻流失的请求记为 $\Delta t$，跨边界未决请求只携带已经历等待时长而不作最终统计。

%%----------- 参考文献 -------------------%%
%在reference.bib文件中填写参考文献，此处自动生成

% 当前正文尚无 \cite 命令，添加正式引用后再启用下行参考文献列表。
% \reference

\end{document}
